\section{Problema detectado y/o faltante}

\subsection{La brecha tecnológica en las asociaciones civiles deportivas}

El análisis de la realidad operativa de los clubes deportivos en la República Argentina revela una profunda asimetría en el acceso a herramientas tecnológicas de gestión y seguridad. Mientras que un puñado de instituciones pertenecientes a la élite de la primera división de fútbol profesional dispone de recursos financieros para adquirir soluciones de \textit{software} a medida e infraestructura de \textit{hardware} propietaria, la inmensa mayoría de las entidades deportivas medianas y pequeñas operan bajo un esquema de marcado rezago tecnológico. Estas instituciones, que forman el tejido social y deportivo del ascenso metropolitano, federal y de las ligas regionales del interior, continúan gestionando sus padrones societarios, finanzas e instalaciones mediante metodologías analógicas o sistemas informáticos fragmentados de baja madurez operativa.

Esta brecha no responde a una falta de voluntad dirigencial, sino a las severas restricciones presupuestarias inherentes a las asociaciones civiles sin fines de lucro en el contexto macroeconómico nacional. Las soluciones comerciales disponibles en el mercado demandan elevados costos de configuración inicial (\textit{setup fees}), licencias de uso restrictivas y la contratación de servicios de consultoría externa para la instalación de servidores locales y su mantenimiento técnico preventivo. Como consecuencia, las comisiones directivas se ven obligadas a prescindir de la digitalización, asumiendo metodologías reactivas basadas en planillas de cálculo de procesamiento manual aislado (como Microsoft Excel) o sistemas de bases de datos locales monolíticos y descentralizados que actúan como meros sistemas \textit{CRUD} básicos (crear, leer, actualizar y borrar). 

La fragmentación y el aislamiento de la información operativa anulan toda posibilidad de implementar auditorías internas eficientes, automatizar la recaudación o consolidar una base de datos única y confiable. Esta carencia estructural no solo genera una alta propensión a errores humanos sistemáticos y pérdidas físicas de registros, sino que expone a las instituciones a severas ineficiencias financieras y vulnerabilidades operativas silenciosas que atentan de manera directa contra su sustentabilidad y saneamiento económico.

\subsection{Caracterización de las tres ineficiencias silenciosas del sector}

El estudio de campo y el relevamiento de requerimientos realizado sobre múltiples clubes del ascenso metropolitano permitieron delimitar de manera formal tres problemáticas críticas de carácter silencioso. Se denominan silenciosas debido a que su impacto financiero e institucional no suele consolidarse en una única métrica evidente, sino que se diluye de forma progresiva a lo largo del ejercicio contable, pasando desapercibido para las auditorías tradicionales hasta que el daño patrimonial es irreversible.

\subsubsection{1. La morosidad societaria y la ineficiencia de cobro transaccional}
El sostenimiento financiero de un club deportivo depende de manera casi exclusiva de la regularidad en el flujo de ingresos por cuotas sociales mensuales y aranceles de disciplinas recreativas. Sin embargo, el sistema tradicional de recaudación presenta una alta fricción operativa que fomenta de forma involuntaria el comportamiento de impago y olvido por parte del asociado:
\begin{itemize}
    \item \textbf{Dependencia de ventanillas físicas}: Gran parte de las instituciones deportivas del interior y de categorías intermedias continúan exigiendo que el socio acuda de forma presencial a la secretaría de la sede social en horarios administrativos restringidos para saldar sus obligaciones de pago.
    \item \textbf{Falta de canales proactivos}: Las secretarías administrativas carecen de mecanismos automatizados para la emisión de alertas de pago personalizadas antes del vencimiento. Los procesos de reclamo de deuda son manuales y reactivos, ejecutándose mediante llamadas telefónicas directas individuales una vez acumulados múltiples meses de mora, momento en el cual la tasa de recuperabilidad del socio decae exponencialmente.
    \item \textbf{Fricción en entornos digitales cerrados}: Las pocas herramientas de autogestión existentes obligan a los usuarios a navegar por portales web complejos, recordar credenciales adicionales o descargar aplicaciones nativas pesadas que no se adaptan a teléfonos móviles de gama baja o con almacenamiento limitado, desincentivando el flujo de pago espontáneo.
\end{itemize}

Esta desconexión crónica entre la tesorería del club y los hábitos transaccionales diarios de los socios consolida una morosidad silenciosa que desfinancia de forma sostenida los presupuestos operativos de las instituciones.

\subsubsection{2. Fraude perimetral en controles de ingreso y préstamo de credenciales}
La seguridad física y el control de accesos a los estadios de fútbol y sedes sociales representan desafíos de alta criticidad operativa, especialmente en el contexto sociocultural de los espectáculos deportivos en Argentina. Los métodos tradicionales de validación perimetral demuestran ser sumamente vulnerables al fraude y la suplantación de identidad:
\begin{itemize}
    \item \textbf{Vulnerabilidad de carnets estáticos}: Las credenciales impresas en tarjetas plásticas de cloruro de polivinilo (\textit{PVC}) con códigos de barras lineales estáticos carecen de cualquier elemento de seguridad dinámica. Esto facilita que los carnets sean prestados de forma sistemática a terceros no autorizados o que se utilicen fotocopias y capturas de pantalla digitales para vulnerar los molinetes perimetrales.
    \item \textbf{Fuga de recaudación neta}: El préstamo masivo de credenciales permite el ingreso de personas morosas o no socias bajo el amparo de un carnet de socio activo regular. Esta duplicidad de accesos no solo constituye una fuga directa de capital por entradas de tribuna no vendidas, sino que distorsiona severamente las métricas de capacidad máxima permitida de los estadios, incrementando el riesgo de clausuras institucionales.
    \item \textbf{Vulnerabilidad en seguridad lógica}: Las dirigencias carecen de un registro nominal confiable en tiempo real de quiénes ingresan de forma efectiva a las tribunas dadores de partido, lo que dificulta la articulación con las fuerzas de seguridad pública en la prevención de desmanes civiles.
\end{itemize}

\subsubsection{3. Ociosidad y desaprovechamiento comercial de activos comunes}
Los clubes de fútbol de tamaño mediano disponen de una gran cantidad de activos compartidos y espacios físicos (tales como canchas de césped sintético auxiliares, salones de usos múltiples, quinchos, gimnasios o sectores de entrenamiento especializado) que permanecen ociosos durante una proporción sustancial de la semana. 

La inexistencia de un canal unificado, digital e inmediato de consulta y reserva impide que el socio autogestione el uso de estas instalaciones. La administración de los turnos se realiza de forma manual en cuadernos de actas físicos por parte de los empleados de la sede, lo que genera duplicidad de reservas, cancelaciones ineficientes y, fundamentalmente, la pérdida de ingresos adicionales que podrían diversificar de manera significativa la matriz de financiamiento de la institución.

\subsection{El desafío físico de infraestructura en entornos masivos}

Cualquier intento de digitalización y segurización del control de accesos mediante herramientas dinámicas tradicionales que dependen de conectividad celular en tiempo real (\textit{online}) se enfrenta a un obstáculo técnico de infraestructura física insalvable en los estadios de fútbol: el colapso sistemático de las redes de comunicación celular en los minutos previos al inicio de un evento deportivo masivo.

La física de la radiofrecuencia y la teoría de colas explican que la aglomeración de decenas de miles de espectadores en un espacio geográfico sumamente delimitado provoca la saturación inmediata de las celdas de telefonía celular móvil de cuarta y quinta generación (4G/5G). Las antenas de los proveedores colapsan ante la alta densidad de conexiones simultáneas. En este escenario de congestión extrema, si la credencial digital del socio requiere de una conexión activa a internet para cargar, actualizarse o realizar peticiones contra un servidor central en la nube, el sistema de ingreso falla de manera inevitable:

\begin{equation}
    \lambda_{\text{peticiones}} \gg \mu_{\text{canal\_comunicación}} \implies \text{Latencia} \to \infty
\end{equation}

Ante este embudo tecnológico, y para evitar desbordes en los accesos de las tribunas por la acumulación de personas en la vía pública, el personal de control se ve obligado sistemáticamente a liberar los ingresos físicos. Esto anula toda efectividad de las medidas de seguridad y posibilita el acceso de personas suspendidas, morosas o con identidades fraudulentas.

SocioUnido resuelve este punto crítico en la experiencia del usuario asegurando que la visualización de la credencial sea completamente independiente de la conectividad de red de la zona del estadio. La aplicación web progresiva (\textit{PWA}) del socio está diseñada para generar los códigos dinámicos de acceso mediante algoritmos de contraseñas de un solo uso basadas en el tiempo (\textit{TOTP}) de forma 100\% fuera de línea (\textit{offline}), permitiendo que el usuario cargue su carnet dinámico de forma inmediata incluso en modo avión, neutralizando de este modo el principal foco de fricción y asegurando que siempre disponga de su credencial al momento de ingresar.

\subsection{La ausencia de alternativas accesibles en el mercado de software}

El relevamiento del estado del arte evidencia un faltante absoluto de plataformas de software que aborden esta problemática de manera integral para los clubes de categorías intermedias. Las soluciones existentes en la industria presentan limitaciones estructurales insalvables para el presupuesto y la infraestructura del ascenso argentino:

\begin{enumerate}
    \item \textbf{Inviabilidad económica}: Los sistemas tradicionales requieren hardware especializado y costoso (lectores fijos RFID, torniquetes electrónicos integrados de alta gama y servidores locales de alto rendimiento) que obligan a las dirigencias a destinar partidas presupuestarias imposibles de amortizar en el mediano plazo.
    \item \textbf{Dependencia de conectividad de red local}: Las pocas alternativas que proponen códigos \textit{QR} variables exigen que el dispositivo lector perimetral mantenga una conexión cableada o inalámbrica a internet de alta velocidad permanente para validar la validez de la firma digital de forma remota, una condición que la infraestructura de los estadios del ascenso no puede garantizar bajo ningún concepto.
    \item \textbf{Complejidad administrativa en interfaces nativas}: Las aplicaciones tradicionales imponen una alta curva de aprendizaje y fricción de uso al obligar a los socios a adaptarse a plataformas complejas y descargar aplicaciones de gran tamaño, ignorando la heterogeneidad tecnológica de la masa social de un club social y deportivo de barrio.
\end{enumerate}

SocioUnido nace como la respuesta de ingeniería de software definitiva a este faltante del mercado. La plataforma se concibe como un ecosistema unificado que democratiza el acceso a la alta tecnología perimetral y analítica mediante una arquitectura elástica, un bot cognitivo accesible sin fricción de descarga y un sistema de acceso inteligente capaz de computar y validar firmas dinámicas de forma offline, neutralizando el colapso de red en eventos masivos y saneando financieramente las instituciones del fútbol argentino.
